\Needspace{9\baselineskip}
\originalchapter{谐振子与Gaussian波包}
\originalsection{升降算符与能级}
在势能极小点附近, Taylor展开的一次项为零, 二次项通常主导, 因而出现 $V(x)=m\omega^2x^2/2$, $\omega>0$. 高次项小才允许此近似; 分子振动和局部束缚常因此近似为谐振子. 定义
\[
a=\sqrt{\frac{m\omega}{2\hbar}}x+\frac{\ii p}{\sqrt{2m\hbar\omega}},
\qquad a^\dagger=\sqrt{\frac{m\omega}{2\hbar}}x-\frac{\ii p}{\sqrt{2m\hbar\omega}}.
\]
在共同的光滑快速下降定义域, $[x,p]=\ii\hbar$给$[a,a^\dagger]=1$, 而 $H=\hbar\omega(a^\dagger a+1/2)$. 令$N=a^\dagger a$, 则对归一态有$\langle N\rangle=\|a\psi\|^2\ge0$. 对易式 $[N,a]=-a$, $[N,a^\dagger]=a^\dagger$ 表示$a$降低、$a^\dagger$升高本征值一个单位.

若$N\ket{\nu}=\nu\ket{\nu}$且$\nu$有限, 连续降阶后的非零态必须保持非负本征值. 降阶范数平方依次含$\nu,\nu-1,\nu-2,\ldots$; 如果$\nu$不是非负整数, 在某一步会出现负范数平方. 因而序列必须终止于$a\ket0=0$, 并有$\nu=n=0,1,\ldots$. 对标准实线谐振子自伴实现, 这些态构成完整基; 完备性是Hermite函数的$L^2$性质, 不由仅有对易式的形式代数自动保证.

基态方程$a\psi_0=0$是 $\psi_0'=-m\omega x\psi_0/\hbar$, 所以
\[
\psi_0(x)=\left(\frac{m\omega}{\pi\hbar}\right)^{1/4}
\e^{-m\omega x^2/(2\hbar)},\qquad E_n=\hbar\omega\left(n+\frac12\right).
\]
由$\ket n=(a^\dagger)^n\ket0/\sqrt{n!}$构造激发态, 可逐次验证$a\ket n=\sqrt n\ket{n-1}$, $a^\dagger\ket n=\sqrt{n+1}\ket{n+1}$. 能量间隔相等源自升降算符代数; 基态能量非零源自$H$中的$1/2$.

\originalsection{位置、动量与期望值}
反解得到 $x=\sqrt{\hbar/(2m\omega)}(a+a^\dagger)$ 和 $p=-\ii\sqrt{m\hbar\omega/2}(a-a^\dagger)$. 在$\ket n$中, 单个升降算符的期望值为零. 展开$x^2$后, $a^2$和$(a^\dagger)^2$也把态移到正交的能级; 只有$a a^\dagger$和$a^\dagger a$留下:
\[
\langle x^2\rangle_n=\frac{\hbar}{m\omega}\left(n+\frac12\right),
\qquad \langle p^2\rangle_n=m\hbar\omega\left(n+\frac12\right).
\]
所以$\Delta x\Delta p=\hbar(n+1/2)$, 仅基态在这些能量本征态中达到$\hbar/2$. 平均动能与势能各为$E_n/2$, 但不是说每次测量都得到这两个平均值.

\originalsection{自由Gaussian的扩散}
给宽度参数$a>0$的初态 $\psi(x,0)=(2\pi a^2)^{-1/4}\e^{-x^2/(4a^2)}$. 自由演化在波数表象中把每个分量乘以$\e^{-\ii\hbar k^2t/(2m)}$. 第4章已算初态波数振幅为$(2a^2/\pi)^{1/4}\e^{-a^2k^2}$. 因而需计算$\int\exp[-(a^2+\ii\hbar t/(2m))k^2+\ii kx]\dd k$. 对$\operatorname{Re}A>0$的Gaussian积分为$\sqrt{\pi/A}\exp[B^2/(4A)]$; 它从实正$A$的积分经解析延拓成立. 取$A=a^2+\ii\hbar t/(2m)$、$B=\ii x$代入. 记$\tau=\hbar t/(2ma^2)$, 得
\[
\psi(x,t)=\frac{(2\pi a^2)^{-1/4}}{\sqrt{1+\ii\tau}}
\exp\left[-\frac{x^2}{4a^2(1+\ii\tau)}\right],
\]
平方根取与$t=0$连续的分支. 取模平方得
\[
\rho(x,t)=\frac1{\sqrt{2\pi}a\sqrt{1+\tau^2}}
\exp\left[-\frac{x^2}{2a^2(1+\tau^2)}\right].
\]
因此$\Delta x(t)=a\sqrt{1+\tau^2}$, 而动量分布只有相位改变, $\Delta p(t)=\hbar/(2a)$保持不变. 这与不同动量分量具有不同速度相符. 初态不带线性相位时中心固定; 若乘以$\e^{\ii p_0x/\hbar}$, 中心变为$p_0t/m$, 宽度保持同样演化.

初态达到不确定性最小值, 后来$\Delta x\Delta p$变大. 量子演化没有违反不确定性关系: 它给下界, 并不要求所有时刻都取等. 演化还产生位置与动量的相关性, 仅用两个独立标准差不能完整描述波包.
